% !TEX program = xelatex
\documentclass[11pt]{ctexart}
%=====================================================================
% theme-violet.tex —— 讲义精讲 · 紫花主题
% 用法：\documentclass[11pt]{ctexart}
%       \input{assets/theme-violet}
%       \begin{document} ... \end{document}
% 只用 ctexart 自带能力 + xcolor/enumitem/fancyhdr/pifont/amsmath，
% 不依赖 tikz/tcolorbox，保证 MiKTeX 可直接编译。
%=====================================================================

\usepackage[a4paper,left=2.0cm,right=2.0cm,top=2.1cm,bottom=2.3cm,%
            headsep=12pt,footskip=22pt]{geometry}
\usepackage{amsmath,amssymb}
\usepackage{xcolor}
\usepackage{array}
\usepackage{enumitem}
\usepackage{pifont}
\usepackage{fancyhdr}
\usepackage[hidelinks]{hyperref}

%---------------------------------------------------------------- 配色
\definecolor{violet}{RGB}{107,78,143}      % 主色：紫
\definecolor{violetdeep}{RGB}{63,44,92}    % 深紫：标题/强调
\definecolor{petal}{RGB}{183,159,208}      % 花瓣：装饰线、花饰
\definecolor{bandviolet}{RGB}{247,243,251} % 浅紫底：档案/提示框
\definecolor{bandcream}{RGB}{252,249,242}  % 米白底：易错/注意框
\definecolor{rulepurple}{RGB}{217,207,230} % 分隔线
\definecolor{inkgray}{RGB}{92,96,110}      % 次要文字

%---------------------------------------------------------------- 字体
% 有 Noto CJK 用 Noto，没有就退回 ctex 默认（Windows 下为宋体/黑体）
\IfFontExistsTF{Noto Serif CJK SC}{\setCJKmainfont{Noto Serif CJK SC}}{}
\IfFontExistsTF{Noto Sans CJK SC}{\setCJKsansfont{Noto Sans CJK SC}}{}

%---------------------------------------------------------------- 版式
\setlist{nosep,leftmargin=1.9em}
\linespread{1.28}
\setlength{\parskip}{3.5pt}
\setlength{\parindent}{0pt}
\setlength{\tabcolsep}{5pt}
\renewcommand{\arraystretch}{1.28}

%---------------------------------------------------------------- 花饰
% \ding{95} = 五瓣小花，\ding{96} = 星形花；都来自 pifont，必定存在
\newcommand{\orn}{\textcolor{petal}{\ding{95}}}          % 行内小标
\newcommand{\ornstar}{\textcolor{petal}{\ding{96}}}      % 星形花
\newcommand{\flowerrule}{%                                % 居中花饰分隔线
  \par\vspace{7pt}%
  \noindent\makebox[\linewidth]{\color{petal}%
    \rule{0.20\linewidth}{0.6pt}\hspace{8pt}{\small\ding{95}}\hspace{4pt}%
    \raisebox{1.1pt}{\tiny\ding{95}}\hspace{4pt}{\small\ding{95}}\hspace{8pt}%
    \rule{0.20\linewidth}{0.6pt}}%
  \par\vspace{7pt}}

%---------------------------------------------------------------- 标题与节头
\newcommand{\doctitle}[3]{%   #1 主标题  #2 副标题  #3 元信息（可用 \\ 分行）
  \begin{center}
    {\color{petal}\small \orn\enspace #2\enspace \ornstar}\\[7pt]
    {\LARGE\bfseries\color{violetdeep} #1}\\[10pt]
    {\color{petal}\rule{0.55\linewidth}{1pt}}\\[4pt]
    {\small\color{inkgray} #3}
  \end{center}\vspace{2pt}}

\newcommand{\hdr}[2]{%   #1 章节号（如 §0）  #2 标题
  \par\vspace{12pt}%
  \noindent\colorbox{violetdeep}{\color{white}\bfseries\small\;#1\;}\hspace{7pt}%
  {\color{violetdeep}\large\bfseries #2}\par\nobreak
  \vspace{3pt}{\color{petal}\hrule height 1.2pt}\vspace{6pt}}

\newcommand{\subhdr}[1]{%
  \par\vspace{6pt}%
  \noindent{\color{violet}\bfseries #1}\par\vspace{1pt}}

%---------------------------------------------------------------- 提示框
\newcommand{\callout}[1]{% 浅紫底：讲义档案、读法说明
  \par\medskip\noindent
  \fcolorbox{rulepurple}{bandviolet}{\parbox{\dimexpr\linewidth-2\fboxsep-2\fboxrule}{\vspace{2pt}#1\vspace{2pt}}}%
  \par\medskip}

\newcommand{\warnbox}[1]{% 米白底：易错、边界条件
  \par\medskip\noindent
  \fcolorbox{rulepurple}{bandcream}{\parbox{\dimexpr\linewidth-2\fboxsep-2\fboxrule}{\vspace{2pt}#1\vspace{2pt}}}%
  \par\medskip}

\newcommand{\gapitem}[2]{% #1 缺口类型（紫） #2 内容
  \par\vspace{2pt}\noindent
  \hangindent=1.4em\hangafter=1
  {\color{petal}\small\ding{96}}\enspace{\bfseries\color{violetdeep}#1}\enspace #2\par}

% 前提旁注：只放在「需要理解的那一步」后面，给陈述 + 本题代入，不重证。
\newcommand{\aside}[1]{%
  \par\vspace{1pt}\noindent
  \hangindent=1.6em\hangafter=1
  {\small\color{inkgray}%
    \hspace*{0.6em}{\color{petal}$\triangleright$}\hspace{2pt}#1}%
  \par\vspace{2pt}}

%---------------------------------------------------------------- 表格
% 本机 MiKTeX 的 colortbl 与 array 版本不匹配（\rowcolor 会报错），
% 故表格统一用「上粗紫线 + 表头 + 中细线」的三线式，不依赖 colortbl。
\newcommand{\tabtop}{\par\vspace{3pt}{\color{violet}\hrule height 1.2pt}\vspace{1pt}}
\newcommand{\tabmid}{\par\vspace{1pt}{\color{petal}\hrule height 0.5pt}\vspace{1pt}}
\newcommand{\tabbot}{\par\vspace{3pt}{\color{rulepurple}\hrule height 0.5pt}\vspace{8pt}}
\newcommand{\thead}[1]{{\bfseries\color{violetdeep}#1}}

%---------------------------------------------------------------- 页眉页脚
\pagestyle{fancy}
\fancyhf{}
\renewcommand{\headrulewidth}{0.4pt}
\renewcommand{\headrule}{\hbox to\headwidth{\color{rulepurple}\leaders\hrule height \headrulewidth\hfill}}
\fancyhead[L]{\small\color{inkgray}\songti \doctag}
\fancyhead[R]{\small\color{inkgray}\songti \coursetag}
\fancyfoot[C]{\small\color{inkgray}\orn\enspace{\color{violet}\thepage}\enspace\orn}

% 每篇讲义在导言区设这两个值
\newcommand{\doctag}{}
\newcommand{\coursetag}{}
\newcommand{\settags}[2]{\renewcommand{\doctag}{#1}\renewcommand{\coursetag}{#2}}

%---------------------------------------------------------------- 数学记号
\newcommand{\vect}[1]{\boldsymbol{#1}}
\newcommand{\dd}{\,\mathrm{d}}
\newcommand{\ddx}{\frac{\mathrm{d}}{\mathrm{d}x}}
\addtolength{\jot}{2pt}



\setlength{\headheight}{22pt}
\usepackage{tikz}
\usetikzlibrary{positioning,arrows.meta,calc}
\usepackage{graphicx}
\usepackage{float}
\usepackage[american]{circuitikz}
\settags{实验 4 · 桥式整流补页}{电工电子技术}
\let\orighdr\hdr
\renewcommand{\hdr}[2]{\needspace{8\baselineskip}\orighdr{#1}{#2}}
\let\origsubhdr\subhdr
\renewcommand{\subhdr}[1]{\needspace{3\baselineskip}\origsubhdr{#1}}

\begin{document}

\doctitle{补页：桥式整流到底怎么走电流}
{配合《实验 4 直流稳压电源》讲义 §2.3；先看图，再读表，最后看波形}
{2026-10-27\ \textperiodcentered\ 建议打印后夹在讲义第 3 页前}

\callout{\textbf{这一页解决什么问题。} 原讲义 §2.3 直接从“四只二极管”跳到“导通表”，中间缺了\textbf{桥的接线图}和\textbf{两条电流路径}：读者看不到 a、b 怎么接、+、− 怎么出、电流从哪只管子进、从哪只管子回，所以“负半周被翻上来”只能死记。下面按“拓扑 \(\to\) 两条路径 \(\to\) 波形”三步补上；看完这三步，再回头看原讲义的表，表就只是图的摘要。}

%=============================================================
\hdr{第一步}{先认桥：四只二极管、两个交流端、两个直流端}

\begin{figure}[H]
\centering
\begin{circuitikz}[scale=0.95,transform shape]
\coordinate (L) at (0,0);
\coordinate (T) at (2.5,2.2);
\coordinate (R) at (5,0);
\coordinate (B) at (2.5,-2.2);
\draw (L) to[D] (T);
\draw (R) to[D] (T);
\draw (B) to[D] (L);
\draw (B) to[D] (R);
\node[font=\small] at (0.65,1.5) {$D_1$};
\node[font=\small] at (4.35,1.5) {$D_2$};
\node[font=\small] at (0.65,-1.5) {$D_3$};
\node[font=\small] at (4.35,-1.5) {$D_4$};
\draw (T) -- (7.0,2.2) to[R, l_=$R_1$] (7.0,-2.2) -- (B);
\draw[<->] (7.7,2.0) -- (7.7,-2.0) node[midway,right] {$U_{o1}$};
\node[left] at (L) {a};
\node[right] at (R) {b};
\node[above] at (T) {+};
\node[below] at (B) {\(-\)};
\node[below,font=\scriptsize] at (2.5,-2.8) {桥的“下顶点”};
\node[above,font=\scriptsize] at (2.5,2.8) {桥的“上顶点”};
\end{circuitikz}
\caption{桥式整流的拓扑：a、b 是交流输入端（接变压器副边），+、− 是直流输出端（接负载 \(R_1\)）。四只二极管两两成对，\(D_1\)、\(D_2\) 的阴极都朝上，\(D_3\)、\(D_4\) 的阳极都朝下。}
\end{figure}

\textbf{先记三个位置，不要急着背导通表：}
\begin{itemize}
  \item \textbf{a、b}：交流输入端。a 是左顶点，b 是右顶点；它们接变压器副边，电压 \(u_2\) 就加在 a、b 之间。
  \item \textbf{+}：上顶点，直流输出的正端；\textbf{−}：下顶点，直流输出的负端。负载 \(R_1\) 接在 +、− 之间。
  \item \textbf{四只管子}：\(D_1\) 在 a 到 + 的边上，\(D_2\) 在 b 到 + 的边上，\(D_3\) 在 − 到 a 的边上，\(D_4\) 在 − 到 b 的边上。注意方向：\(D_1\)、\(D_2\) 的箭头都指向 +；\(D_3\)、\(D_4\) 的箭头都从 − 指向 a/b。
\end{itemize}

%=============================================================
\clearpage
\hdr{第二步}{分两个半周看：哪两只管子导通、电流怎么走}

\begin{figure}[H]
\centering
\begin{minipage}[t]{0.48\linewidth}
\centering
\begin{circuitikz}[scale=0.70,transform shape]
\coordinate (L) at (0,0); \coordinate (T) at (2.5,2.2);
\coordinate (R) at (5,0); \coordinate (B) at (2.5,-2.2);
% gray bridge
\draw[gray] (L) to[D] (T);
\draw[gray] (R) to[D] (T);
\draw[gray] (B) to[D] (L);
\draw[gray] (B) to[D] (R);
% red path
\draw[red,very thick] (L) to[D] (T) -- (6.6,2.2) to[R] (6.6,-2.2) -- (B) to[D] (R);
\draw[red,very thick,->] (5.8,1.6) -- (5.8,0.4);
\node[left,font=\small] at (L) {a\,+};
\node[right,font=\small] at (R) {b\,\(-\)};
\node[above,font=\small] at (T) {+};
\node[below,font=\small] at (B) {\(-\)};
\node[font=\scriptsize,red] at (5.8,-1.1) {$I_L$};
\node[font=\small] at (0.65,1.5) {$D_1$};
\node[font=\small] at (4.35,1.5) {$D_2$};
\node[font=\small] at (0.65,-1.5) {$D_3$};
\node[font=\small] at (4.35,-1.5) {$D_4$};
\end{circuitikz}
\par\small (a) \(u_2>0\)：a 正、b 负；\(D_1\)、\(D_4\) 导通（红），\(D_2\)、\(D_3\) 截止（灰）。
\end{minipage}\hfill
\begin{minipage}[t]{0.48\linewidth}
\centering
\begin{circuitikz}[scale=0.70,transform shape]
\coordinate (L) at (0,0); \coordinate (T) at (2.5,2.2);
\coordinate (R) at (5,0); \coordinate (B) at (2.5,-2.2);
\draw[gray] (L) to[D] (T);
\draw[gray] (R) to[D] (T);
\draw[gray] (B) to[D] (L);
\draw[gray] (B) to[D] (R);
% blue path
\draw[blue,very thick] (R) to[D] (T) -- (6.6,2.2) to[R] (6.6,-2.2) -- (B) to[D] (L);
\draw[blue,very thick,->] (5.8,1.6) -- (5.8,0.4);
\node[left,font=\small] at (L) {a\,\(-\)};
\node[right,font=\small] at (R) {b\,+};
\node[above,font=\small] at (T) {+};
\node[below,font=\small] at (B) {\(-\)};
\node[font=\scriptsize,blue] at (5.8,-1.1) {$I_L$};
\node[font=\small] at (0.65,1.5) {$D_1$};
\node[font=\small] at (4.35,1.5) {$D_2$};
\node[font=\small] at (0.65,-1.5) {$D_3$};
\node[font=\small] at (4.35,-1.5) {$D_4$};
\end{circuitikz}
\par\small (b) \(u_2<0\)：b 正、a 负；\(D_2\)、\(D_3\) 导通（蓝），\(D_1\)、\(D_4\) 截止（灰）。
\end{minipage}
\caption{两条路径：红色是正半周，蓝色是负半周。两次都从 + 流出、从 − 流回，所以 \(R_1\) 上的电流方向始终从上到下。}
\end{figure}

\textbf{怎么读这两张图（按顺序做）：}
\begin{enumerate}
  \item 先看 \(u_2\) 的极性：a 正还是 b 正。阳极接高电位、阴极接低电位的那只二极管才会导通；反偏的截止。
  \item 从 a（或 b）出发，沿导通二极管走到 +，再穿过 \(R_1\) 到 −，最后沿另一只导通二极管回到 b（或 a）。
  \item 确认 \(I_L\) 的方向：两张图里 \(I_L\) 都从上往下。方向不变，大小仍在脉动，这就是“全波脉动直流”。
\end{enumerate}

\warnbox{\textbf{一句话记住：}正半周走 \(D_1\to R_1\to D_4\)，负半周走 \(D_2\to R_1\to D_3\)。两只管子一组，轮流把 a、b 接到 +、− 上；负载始终接在 +、− 之间，所以电流不反向。}

\flowerrule
\callout{\textbf{回到原讲义。} 现在再翻到 §2.3 的导通表：它其实就是把上面两张小图压缩成了两行。看到“\(u_2>0\)：\(D_1,D_4\)”时，脑子里应该浮现红色路径；看到“\(u_2<0\)：\(D_2,D_3\)”时，应该浮现蓝色路径。如果还能浮现不出来，就自己拿一张纸，把图 1 重画一遍，再用红、蓝笔各描一次路径——这比反复读表有效。}

%=============================================================
\hdr{第三步}{再看波形：负半周被“翻”到正半周}

\begin{figure}[H]
\centering
\begin{tikzpicture}[scale=0.9,>=Latex]
% u2
\begin{scope}[xshift=0cm]
\draw[->] (0,0) -- (7.2,0) node[right] {\scriptsize $t$};
\draw[->] (0,-1.8) -- (0,1.9) node[above] {\scriptsize $u_2$};
\draw[thick,blue,domain=0:720,samples=240] plot ({\x/100},{1.5*sin(\x)});
\draw[dashed] (0,0) -- (7.2,0);
\node[font=\scriptsize] at (1.8,1.65) {正半周};
\node[font=\scriptsize] at (5.4,-1.65) {负半周};
\foreach \x in {1.8,3.6,5.4} \draw (\x,0.06) -- (\x,-0.06);
\node[below,font=\scriptsize] at (1.8,0) {$\pi$};
\node[below,font=\scriptsize] at (3.6,0) {$2\pi$};
\node[below,font=\scriptsize] at (5.4,0) {$3\pi$};
\node[below,font=\scriptsize] at (0,0) {0};
\end{scope}
% uo1
\begin{scope}[xshift=0cm,yshift=-4.6cm]
\draw[->] (0,0) -- (7.2,0) node[right] {\scriptsize $t$};
\draw[->] (0,-0.2) -- (0,1.9) node[above] {\scriptsize $u_{o1}$};
\draw[thick,red,domain=0:720,samples=240] plot ({\x/100},{1.5*abs(sin(\x))});
\draw[dashed] (0,0.955) -- (7.0,0.955);
\node[font=\scriptsize] at (2.8,1.15) {$0.9U_2$};
\foreach \x in {1.8,3.6,5.4} \draw (\x,0.06) -- (\x,-0.06);
\node[below,font=\scriptsize] at (1.8,0) {$\pi$};
\node[below,font=\scriptsize] at (3.6,0) {$2\pi$};
\node[below,font=\scriptsize] at (5.4,0) {$3\pi$};
\node[below,font=\scriptsize] at (0,0) {0};
\node[font=\scriptsize] at (3.6,1.7) {周期变成 $T/2=10\,\mathrm{ms}$};
\end{scope}
% arrow between panels
\draw[->,thick] (7.6,0.2) -- (7.6,-4.4);
\node[font=\scriptsize,right] at (7.7,-2.1) {负半周翻上来};
\end{tikzpicture}
\caption{上：副边交流 \(u_2\) 有正有负；下：桥式整流后 \(u_{o1}=|u_2|\)，负半周被翻到正半周，脉动周期变成 \(T/2=10\,\mathrm{ms}\)，平均值约 \(0.9U_2\)。}
\end{figure}

\textbf{波形和路径的对应：}
\begin{itemize}
  \item 上图的每一个正半周，对应 (a) 的红路：\(D_1\to R_1\to D_4\)；
  \item 上图的每一个负半周，对应 (b) 的蓝路：\(D_2\to R_1\to D_3\)；
  \item 两条路都让 \(R_1\) 的电流从上往下，所以下图全是“正拱”，没有负半周；这就是“全波”。
\end{itemize}

%=============================================================
\hdr{自检与常见误解}

\subhdr{三个自检问题（能答出来，就真的看懂了）}

\begin{enumerate}
  \item 正半周（a 正、b 负）时，电流从 a 出发，经过哪两只二极管回到 b？——答：\(D_1\to R_1\to D_4\)。
  \item 负半周（b 正、a 负）时，电流从 b 出发，经过哪两只二极管回到 a？——答：\(D_2\to R_1\to D_3\)。
  \item 如果把 \(D_1\) 反接，正半周会发生什么？——答：正半周没有通路，输出缺一半，波形变成“半波”，平均值明显下降；严重时还可能短路。这个反例能帮你记住 \(D_1\) 的方向。
\end{enumerate}

\subhdr{四个常见误解}

\begin{center}
\small
\begin{tabular}{@{}p{6.2cm}p{9.4cm}@{}}
\hline
\thead{误解} & \thead{纠正} \\
\hline
“四只二极管同时导通” & 每个半周只有两只导通；另外两只反偏截止。 \\
“桥式整流是半波” & 桥式是全波：正、负半周都用上，脉动频率 \(100\,\mathrm{Hz}\)。 \\
“输出已经是直流” & 输出是\textbf{脉动直流}，方向不变但大小仍在变；还要滤波。 \\
“平均值是 \(0.45U_2\)” & \(0.45U_2\) 是半波；桥式全波是 \(0.9U_2\)。 \\
\hline
\end{tabular}
\end{center}

\end{document}
